\src{172}
$\symup{0\longrightarrow K'\longrightarrow K\longrightarrow K''\longrightarrow 0}$, telle que $\symup{K'^{i}}$ s'identifie à un facteur direct de $\symup{K^{i}}$ pour tout $\symup{i}$,
donne naissance à une suite exacte pour les complexes de groupes
$\symup{\operatorname{Hom}(K'',L)}$, etc., d'où un opérateur cobord
$\symup{H^{i}(K',L)\longrightarrow H^{i+1}(K'',L)}$.
On définit de même les opérateurs cobord relatifs à une suite exacte en $\symup{L}$. Les
accouplements (2.2) sont compatibles, au sens usuel, avec ces opérateurs cobord.

Supposons maintenant que $\symup{C}$ soit une catégorie telle que tout élément $\symup{A}$
de $\symup{C}$ admette une résolution injective $\symup{C(A)}$. On constate alors, utilisant
une des nombreuses variantes du théorème du bicomplexe, que
\[
  \symup{H^{*}\bigl(C(A),C(B)\bigr)
  =H^{*}\bigl(\operatorname{Hom}(C(A),C(B))\bigr)}
\]
est canoniquement isomorphe à
\[
  \symup{H^{*}\bigl(\operatorname{Hom}(A,C(B))\bigr)
  =\operatorname{Ext}^{*}(A,B).}
\]
Les opérateurs cobord envisagés plus haut donnent les homomorphismes cobord des
$\symup{\operatorname{Ext}}$. De plus, les accouplements (2.2) donnent ici des
accouplements associatifs, compatibles avec les homomorphismes cobord :
\begin{equation}\tag{2.3}
  \symup{\operatorname{Ext}^{*}(A,B)\times\operatorname{Ext}^{*}(B,C)
  \longrightarrow\operatorname{Ext}^{*}(A,C)}
\end{equation}
En particulier, $\symup{\operatorname{Ext}^{*}(A,A)}$ est un anneau gradué associatif
avec unité, etc. (On montre de façon analogue que les foncteurs
$\symup{\operatorname{Ext}}$ opèrent dans les foncteurs dérivés d'un foncteur
quelconque ; nous ne nous servirons pas ici de ce fait.)

Dans le cas où la catégorie envisagée est la catégorie $\symup{\Ccat^{\Osh}}$ des
$\symup{\Osh}$-Modules sur $\symup{X}$, on obtient donc des accouplements
\begin{equation}\tag{2.4}
  \symup{\operatorname{Ext}_{\Osh}^{p}(X;A,B)
  \times\operatorname{Ext}_{\Osh}^{q}(X;B,C)
  \longrightarrow\operatorname{Ext}_{\Osh}^{p+q}(X;A,C)}
\end{equation}
qui peuvent s'expliciter comme il a été dit. D'ailleurs, les mêmes
développements, mais où on remplace la catégorie des groupes abéliens par la
catégorie des faisceaux abéliens sur $\symup{X}$, et les foncteurs $\symup{\operatorname{Hom}}$
par les foncteurs $\symup{\Homsh}$, définissent aussi des accouplements, ayant les mêmes
propriétés formelles, et de \fquote{nature locale} cette fois-ci :
\begin{equation}\tag{2.5}
  \symup{\Extsh_{\Osh}^{p}(A,B)\times\Extsh_{\Osh}^{q}(B,C)
  \longrightarrow\Extsh_{\Osh}^{p+q}(A,C)}
\end{equation}
Ces derniers se précisent si on remarque que les homomorphismes (1.8) sont
compatibles avec les accouplements entre les $\symup{\operatorname{Ext}}$.

